\section{Qualitative Error Taxonomy and Case Studies}
\label{sec:error_taxonomy}

\subsection{Fine-Grained Bengali Hallucination Taxonomy}
To analyze generative failure modes beyond binary factuality, Matibhrom adopts an eight-category taxonomy (Codebook v1.0.0) synthesized from classic hallucination frameworks~\citep{maynez2020faithfulness,ji2023survey} and tailored to Bengali linguistic structures (Table~\ref{tab:codebook_taxonomy}).

\begin{table*}[t]
\centering
\small
\begin{tabular}{lp{4.2cm}p{9.0cm}}
\toprule
Code & Category Name & Operational Definition per Bengali Linguistic Annotation Codebook \\
\midrule
\textbf{H1} & Factual Fabrication / Contradiction & Directly contradicts the gold answer or asserts empirical falsehoods (e.g., historical dates, sovereign state boundaries). \\
\textbf{H2} & Extrinsic Unsupported Claim & Asserts factual claims that cannot be verified or substantiated by either parametric knowledge or retrieved passages. \\
\textbf{H3} & Fabricated Citation / Entity & Confabulates non-existent institutional bodies, fake statutory code sections, or invented academic/legal citations. \\
\textbf{H4} & Numerical / Temporal Distortion & Fails on numerical quantities, constitutional article numbers, clinical medication dosages, or calendar dates. \\
\textbf{H5} & Erroneous Deductive Reasoning & Accepts false premises introduced by adversarial prompts or draws invalid deductive conclusions from evidence. \\
\midrule
\textbf{A}  & Appropriate Selective Abstention & Correctly identifies evidence insufficiency under condition $C_3$ or unanswerable items and issues standardized refusal. \\
\textbf{R}  & Unwarranted Refusal & Erroneously refuses to answer despite the presence of unambiguous gold evidence in the provided context ($C_1$). \\
\textbf{none}& Fully Faithful Response & Contains no factual errors, contradictions, or ungrounded additions; strictly faithful to evidence. \\
\bottomrule
\end{tabular}
\caption{Matibhrom Fine-Grained Error Taxonomy (Codebook v1.0.0).}
\label{tab:codebook_taxonomy}
\end{table*}

\subsection{Qualitative Case Studies}
To elucidate the concrete linguistic mechanisms underlying quantitative errors, we analyze three representative qualitative case studies observed during benchmark evaluation.

\subsubsection{Case Study 1: Cross-Jurisprudential Transposition (Categories H1 and H3)}
In this query, the model is prompted in native Standard Bangla: \foreignlanguage{bengali}{"বাংলাদেশের সংবিধান কবে গৃহীত হয়?"} (\emph{"When was the Constitution of Bangladesh adopted?"}), for which the canonical ground truth is \foreignlanguage{bengali}{"৪ নভেম্বর ১৯৭২"} (4 November 1972). Under unconstrained dense retrieval ($C_2$), the retriever surfaces topical distractors discussing local administrative government structures (\texttt{GENERAL-REF-344}). Unable to identify the target date in the context, Qwen 1.7B generates:
\begin{quote}
    \foreignlanguage{bengali}{\emph{"ভারতের সংবিধান ২৬ নভেম্বর ১৯৪৯ সালে গণপরিষদে গৃহীত হয়েছিল এবং ২৬ জানুয়ারি ১৯৫০ সালে কার্যকর হয়। অনুরূপভাবে বাংলাদেশের সংবিধানের ক্ষেত্রেও গণপরিষদীয় প্রক্রিয়া অনুসৃত হয়।"}} \\
    \emph{("The Constitution of India was adopted in the Constituent Assembly on 26 November 1949 and took effect on 26 January 1950. Similarly, in the case of the Bangladesh Constitution, a constituent assembly process was followed.")}
\end{quote}
Because pretraining corpora for Bengali are orders of magnitude smaller than English or Hindi corpora, the model's parametric prior for \foreignlanguage{bengali}{"সংবিধান গৃহীত"} (\emph{"constitution adopted"}) is dominated by Indian constitutional history. When the retrieved context lacks the exact Bangladeshi date, the model suffers from \emph{cross-jurisprudential transposition}, substituting the prompt's sovereign entity with high-resource regional priors and asserting a false historical date.

\subsubsection{Case Study 2: Abstention Rectification under Scaling ($C_2 \rightarrow C_3$)}
The second case examines the prompt \foreignlanguage{bengali}{"বাংলাদেশের সংবিধান অনুযায়ী রাষ্ট্রের সব ক্ষমতার মালিক কারা?"} (\emph{"According to the Constitution of Bangladesh, who is the owner of all power in the Republic?"}), targeting Article 7(1) with gold answer \foreignlanguage{bengali}{"জনগণ"} (\emph{"The People"}). The retrieved context under conditions $C_2$ and $C_3$ consists of administrative transport passages discussing deep-sea port developments. Under unconstrained RAG ($C_2$), Gemma 4B attempts to bridge the semantic gap between the query and the off-target evidence, confabulating:
\begin{quote}
    \foreignlanguage{bengali}{\emph{"বাংলাদেশের সমুদ্রবন্দর ও আন্তর্জাতিক বাণিজ্যের জন্য গভীর সমুদ্রবন্দর হিসেবে মাতারবাড়ী বন্দর প্রকল্প রাষ্ট্রের অর্থনৈতিক ক্ষমতার মূল কেন্দ্র হিসেবে কাজ করে।"}} \\
    \emph{("As a deep sea port for Bangladesh's maritime and international commerce, the Matarbari port project serves as the core center of the state's economic power.")}
\end{quote}
Here, the model forces an artificial semantic connection between \foreignlanguage{bengali}{"ক্ষমতার মালিক"} (\emph{"owner of power"}) and the passage's mention of \foreignlanguage{bengali}{"অর্থনৈতিক ক্ষমতা"} (\emph{"economic power"}). However, when augmented with the selective abstention instruction under condition $C_3$, Gemma 4B correctly recognizes that Article 7(1) is absent from the context and outputs the standardized refusal:
\begin{quote}
    \foreignlanguage{bengali}{\emph{"প্রমাণ অপর্যাপ্ত, তাই আমি উত্তর দিতে পারছি না।"}} \\
    \emph{("Evidence is insufficient, therefore I cannot answer.")}
\end{quote}
This illustrates appropriate selective abstention (Category A), wherein parameter capacity enables the model to suppress confabulation in the presence of distracting information.

\subsubsection{Case Study 3: Banglish Orthographic Dispersion and Token Loops (Category H4)}
The third case evaluates informal Latin transliteration using the Banglish prompt: \emph{"bangladesh er shongbidhan kobe pass hoyesilo?"}. In closed-book generation ($C_0$), Qwen 1.7B generates:
\begin{quote}
    \emph{"Bangladesh er shongbidhan 1971 shaler 26 March e pass hoy ebong shongbidhan board er maddhome shongbidhan shongbidhan shongbidhan protishthito hoy."} \\
    \emph{("The Constitution of Bangladesh was passed on 26 March 1971 and through the constitution board the constitution constitution constitution was established.")}
\end{quote}
In Latin transliteration, the subword tokenizer decomposes \emph{"shongbidhan"} into four sub-character byte tokens (\texttt{sh}, \texttt{ong}, \texttt{bid}, \texttt{han}), compared to a single morphemic unit in native script. The elevated entropy across subword probability distributions causes the autoregressive decoder to enter a degenerate repetition loop, confusing the Independence Declaration date (26 March 1971) with constitutional adoption (4 November 1972) and inventing a fictional \emph{"constitution board"}.
